\chapter{Introduction} % Or \chapter{Introduction}
\label{chap:introduction}


Magnetic resonance imaging (MRI) has consolidated its role as a powerful, non-invasive tool for visualizing the complex anatomy of the human eye and surrounding orbital structures \cite{townsend_clinical_2008, fanea_review_2012}. Its applications in ophthalmology continue to expand, offering crucial insights into the diagnosis, monitoring, and understanding of a wide range of ocular pathologies \cite{duong_magnetic_2014, niendorf_ophthalmic_2021}. The advent of MRI techniques promises to offer even finer anatomical details. This additional precision is essential for the assessment of subtle structural changes and for advancing research in ocular physiology and pathophysiology, including functional studies such as retinal fMRI and vascular oxygenation research \cite{duong_functional_2002, zhang_magnetic_2011}.

However, the diagnostic and experimental utility of ocular MRI relies critically on obtaining high-quality images. The unique characteristics of the eye (small size, mobility, and complex layered structures) present significant imaging challenges \cite{fanea_review_2012, duong_magnetic_2014}. Image quality can be degraded by various artifacts, including those related to involuntary saccades and blinks, as well as physiological eye movements \cite{franceschiello_3-dimensional_2020}. Further complications arise from the inherent complexities of MRI acquisition, such as susceptibility artifacts near air-tissue interfaces and tradeoffs between spatial resolution, signal-to-noise ratio (SNR), and acquisition time \cite{niendorf_ophthalmic_2021}. These factors can obscure critical diagnostic information, lead to misinterpretation, or reduce the reliability of quantitative measurements \cite{schmidt_association_2019}. Therefore, quality control (QC) methods for ocular MRI (MR-Eye) data are essential to ensure data fidelity and reliability.
Although automated quality control frameworks have seen considerable development and adoption for other anatomical regions, including neuroimaging (MRIQC for brain MRI \cite{esteban_mriqc_2017}, which forms the basis of FetMRQC \cite{noauthor_medical-image-analysis-laboratoryfetmrqc_sr_nodate}), they are often not directly transferable or optimal for ocular imaging due to the eyes specific anatomical features, distinct image contrast properties, and common artifact profiles. Manual quality control, involving visual inspection by trained experts, remains a common practice. However, manual assessment is inherently subjective, labor-intensive, and subject to inter- and intra-rater variability, making it a bottleneck for large-scale clinical studies, longitudinal research, and standardized multicenter trials.

This project, MREyeQC, directly addresses this gap. It aims to improve and adapt the FetMRQC framework, originally designed for fetal brain MRI quality control \cite{noauthor_medical-image-analysis-laboratoryfetmrqc_sr_nodate}, to establish a robust and automated quality control pipeline specifically tailored for MR-Eye. The main objective is to develop a system capable of quantitatively assessing the quality of ocular MR images through the calculation of a relevant suite of image quality measures (IQMs). These IQMs are designed to encompass a wide range of quality aspects, including general image properties (e.g., noise, sharpness, contrast), morphological features of key ocular structures (such as the lens and globe), and indicators of common imaging artifacts. Subsequently, these quantitative measurements are used to train machine learning models aimed at predicting image quality with accuracy and consistency comparable to, or even superior to, that of human experts.
